\documentclass[11pt]{article}
\usepackage[margin=1in]{geometry}
\usepackage{amsmath,amssymb,amsthm}

\theoremstyle{plain}
\newtheorem{theorem}{Theorem}
\newtheorem{lemma}[theorem]{Lemma}

\begin{document}
\pagestyle{empty}

\begin{lemma}[Cycle decomposition of a closed walk]
\label{lem:decomp}
Let $D$ be a digraph and let $W = w_0, w_1, \dots, w_L$ (with $w_0 = w_L$) be a closed
directed walk of length $L$ (i.e.\ $w_{t-1} \to w_t$ is an arc of $D$ for each
$t = 1, \dots, L$). Then there exist directed cycles $C_1, \dots, C_m$ of $D$ (not
necessarily distinct, possibly $m=0$) whose lengths satisfy
$|C_1| + \cdots + |C_m| = L$.
\end{lemma}

\begin{proof}
Strong induction on $L$. If $L = 0$, take $m=0$; the empty sum is $0 = L$.

If $L \ge 1$: either $w_0, w_1, \dots, w_{L-1}$ are all distinct, in which case $W$
itself is already a simple directed cycle of length $L$ (take $m=1$, $C_1 = W$); or
some vertex repeats among them, i.e.\ there exist $0 \le a < b \le L-1$ with
$w_a = w_b$. In the latter case, split $W$ into two shorter closed walks:
\[
W_1 = w_a, w_{a+1}, \dots, w_b \quad (\text{length } b-a \ge 1, \text{ closed since } w_a = w_b),
\]
\[
W_2 = w_0, \dots, w_a, w_{b+1}, \dots, w_L \quad (\text{length } L - (b-a), \text{ closed since } w_0 = w_L).
\]
Both $W_1$ and $W_2$ are closed directed walks of length strictly less than $L$
(since $b - a \ge 1$), so the inductive hypothesis applies to each, giving cycles
decomposing $W_1$ and $W_2$ separately. Their union decomposes $W$, with total length
$(b-a) + (L-(b-a)) = L$.
\end{proof}

\begin{theorem}[B]
If $D$ is strongly connected, every directed cycle of $D$ has even length, and
$\delta^+(D) \ge 2$, then $D$ has two disjoint 2-d kernels.
\end{theorem}

\begin{proof}
Fix any vertex $r$. \emph{Claim: every directed path from $r$ to a given vertex $v$
has the same length parity, regardless of which path is taken.} Suppose $P$ and $Q$
are two $r \to v$ paths with $|P| \not\equiv |Q| \pmod 2$. Since $D$ is strongly
connected, fix a directed path $R$ from $v$ back to $r$. Then $P$ followed by $R$,
and $Q$ followed by $R$, are each closed directed walks based at $r$, of lengths
$|P|+|R|$ and $|Q|+|R|$ respectively. By Lemma~\ref{lem:decomp}, each decomposes into
directed cycles of $D$; since every directed cycle of $D$ has even length by
hypothesis, each decomposition's total length is a sum of even numbers, hence even.
So $|P|+|R|$ and $|Q|+|R|$ are both even, forcing
\[
|P| \equiv -|R| \equiv |Q| \pmod 2,
\]
contradicting $|P| \not\equiv |Q|$. So the claim holds: every vertex $v$ has a
well-defined parity $\mathrm{par}(v) \in \{0,1\}$, the length mod 2 of any directed
path from $r$ to $v$ (taking $\mathrm{par}(r) = 0$ via the length-0 path).

Let $V_0 = \{v : \mathrm{par}(v) = 0\}$, $V_1 = \{v : \mathrm{par}(v) = 1\}$; these
partition $V$ since $D$ is strongly connected (every vertex is reachable from $r$).
\emph{Every arc flips parity}: if $u \to v$ is an arc, take any $r \to u$ path $P$ and
extend it by this arc to an $r \to v$ path of length $|P|+1$; by the claim, \emph{every}
$r \to v$ path — including this one — has parity $|P|+1 \bmod 2$, so
$\mathrm{par}(v) = 1 - \mathrm{par}(u)$.

Two consequences follow. First, no arc can join two vertices of equal parity (that
would contradict the flip just shown), so $V_0$ and $V_1$ are each independent.
Second, for $v \in V_i$, every out-neighbour of $v$ lies in $V_{1-i}$. So setting
$S = V_{1-i}$: every vertex outside $S$ (i.e.\ every vertex of $V_i$) has \emph{all}
of $N^+(v)$ contained in $S$, and $|N^+(v)| = d^+(v) \ge \delta^+(D) \ge 2$. Hence
$|N^+(v) \cap S| \ge 2$ for every such $v$, so $S = V_{1-i}$ is a 2-d kernel — for each
$i \in \{0,1\}$. This gives two 2-d kernels, $V_0$ and $V_1$, disjoint since they
partition $V$.

(Both are non-empty: $r \in V_0$ via the length-0 path, and $r$ has out-degree
$\ge 2$, so at least one out-neighbour of $r$ lies in $V_1$.)
\end{proof}

\end{document}
